\justifying
\NSFCBodyText

本次拟申报题目为“多源气象驱动的新疆光伏功率多尺度概率预测与爬坡风险预警研究”，按中国沙漠气象科学研究基金自由申报类组织研究，执行期2027年1月至2028年12月，建议申请3万元。当前文本为借用NSFC结构形成的正文工作稿，正式申报须转入本基金官方申请书，并完成研究摘要、经费预算、成员分工、汇总表及签章材料。

已核验的2026年度指南将未结题基本科研业务费等项目列入限项，并规定截止日前提交结题验收材料的相应例外。申请人主持项目XJEDU2026P005的具体适用解释及截止日前状态，以\textbf{【待补 F-XJ-05：基金管理方的限项解释与相关验收材料状态】}为据；不能仅因项目来自高校、名称不同或研究内容有区分而自行认定不受限。依托单位的未验收项目情况、其他限制情形及本轮真实性承诺，应根据\textbf{【待补 F-XJ-08：申请人及依托单位的本轮核对与签章确认材料】}如实完成。

基金管理办法要求成果第一归属单位为中国气象局乌鲁木齐沙漠气象研究所，论文第一资助基金列为中国沙漠气象科学研究基金并注明获批编号，相关技术档案按规定移交。具体落实以\textbf{【待补 F-XJ-06：依托单位对成果归属、基金标注与技术档案移交的确认】}为据，本文不代替申请人或单位接受协议。数据、代码与知识产权的具体许可须与上述要求相容；未经权利人许可的原始场站数据不能在成果材料中擅自公开。

本轮尚未取得资助批准号，不填写拟造编号；尚未提交官方表格、签章材料或电子邮件。已确认的研究基础与拟开展任务分开列报，待补事实和设计性暂定值补齐后，应重新核查资格、引用、篇幅及正式附件。
